\part{Querschnittsdaten}

\chapter{Momente statt Verteilungen}

\section{Der gemeinsame Nenner}

\begin{merksatz}
\boxtitel{Drei Schätzer, ein Prinzip}
\begin{center}
{\small
\renewcommand{\arraystretch}{1.55}
\begin{tabular}{@{}P{2.8cm}P{5.8cm}@{}}
\toprule
\textbf{Kleinste Quadrate} & $\mathbb{E}[x\,(y-\beta x)] = 0$\\
\textbf{Instrument} & $\mathbb{E}[z\,(y-\beta x)] = 0$\\
\textbf{Mittelwert} & $\mathbb{E}[y - \mu] = 0$\\
\bottomrule
\end{tabular}}
\end{center}

\begin{center}
\begin{minipage}{0.9\textwidth}
\textbf{Jeder dieser Schätzer ist die Lösung einer Gleichung, die sagt: Ein bestimmter Erwartungswert ist null.}
\end{minipage}
\end{center}

\begin{center}
\begin{minipage}{0.88\textwidth}
\kursiv{\textbf{Die Maximum-Likelihood-Schätzung aus Band 17 verlangt die ganze Verteilung}; hier verlangt man nur, dass bestimmte Mittelwerte stimmen. Das ist weniger --- und deshalb robuster. Ökonomische Theorie liefert fast nie Verteilungen, aber oft Bedingungen der Form \enquote{im Mittel gleicht sich das aus}: Gleichgewicht, Erwartungstreue, Fairness. Jede solche Bedingung ist ein Moment.}
\end{minipage}
\end{center}
\end{merksatz}

\begin{definition}
\boxtitelm{Definition 4.1 --- Momentenbedingung und Momentenschätzer}
Eine \Begriff{Momentenbedingung} ist eine Gleichung
\[ \mathbb{E}\bigl[g(w,\theta)\bigr] = 0 , \]
in der $w$ die beobachteten Größen und $\theta$ den Parameter bezeichnet. Der \Begriff{Momentenschätzer} ersetzt den Erwartungswert durch das Stichprobenmittel und löst
\[ \frac1n\sum_{i=1}^{n} g(w_i,\hat\theta) = 0 . \]
\end{definition}

\begin{merksatz}
\boxtitel{Satz 4.2 --- Konsistenz}
Gilt die Momentenbedingung, ist $\theta$ dadurch eindeutig bestimmt und ist $g$ stetig, so strebt $\hat\theta$ gegen $\theta$.

\schritt{Der Gedanke}
Nach dem Gesetz der großen Zahlen strebt das Stichprobenmittel gegen den Erwartungswert --- für jedes feste $\theta$. Die Lösung einer Gleichung, deren linke Seite gleichmäßig gegen eine Grenzfunktion strebt, strebt gegen die Lösung der Grenzgleichung, sofern diese eindeutig ist. $\blacksquare$

\begin{center}
\begin{minipage}{0.88\textwidth}
\kursiv{\textbf{Die Eindeutigkeit ist die Identifikation aus Kapitel~2} --- hier in der Form: Die Momentenbedingung hat genau eine Lösung. Die Gleichmäßigkeit der Konvergenz ist technisch und wird meist stillschweigend vorausgesetzt; wir übergehen ihren Nachweis --- die zweite Auslassung dieses Bandes.}
\end{minipage}
\end{center}
\end{merksatz}

\section{Mehr Bedingungen als Parameter}

\begin{merksatz}
\boxtitel{Das Problem}
Mit zwei Instrumenten $z_1$, $z_2$ und einem Parameter $\beta$ hat man zwei Gleichungen für eine Unbekannte:
\[ \frac1n\sum z_{1i}(y_i-\beta x_i) = 0, \qquad \frac1n\sum z_{2i}(y_i-\beta x_i) = 0 . \]

\begin{center}
\begin{minipage}{0.9\textwidth}
\textbf{In der Stichprobe lassen sich beide nie gleichzeitig exakt erfüllen. Man muss einen Kompromiss wählen.}
\end{minipage}
\end{center}
\end{merksatz}

\begin{definition}
\boxtitelm{Definition 4.3 --- verallgemeinerte Momentenmethode}
Sei $\bar g(\theta) = \frac1n\sum g(w_i,\theta)$ der Vektor der Stichprobenmomente und $W$ eine positiv definite Gewichtsmatrix. Der \Begriff{GMM-Schätzer} minimiert
\[ Q(\theta) = \bar g(\theta)^{\top}\,W\,\bar g(\theta) . \]
\end{definition}

\begin{merksatz}
\boxtitel{Satz 4.4 --- die beste Gewichtung}
Unter allen Gewichtsmatrizen liefert
\[ W^{*} = S^{-1}, \qquad S = \operatorname{Var}\bigl[g(w,\theta)\bigr] \]
den Schätzer mit kleinster Varianz.

\schritt{Der Gedanke}
Momente, die stark streuen, sind weniger verlässlich und sollen weniger zählen; Momente, die miteinander korrelieren, tragen gemeinsam weniger bei, als ihre Zahl vermuten lässt. Beides leistet die inverse Varianzmatrix --- derselbe Gedanke wie bei den gewichteten kleinsten Quadraten in Band 17, Kapitel~8. $\blacksquare$

\begin{center}
\begin{minipage}{0.88\textwidth}
\kursiv{\textbf{In der Praxis rechnet man zweimal}: erst mit einer beliebigen Gewichtung, um $\theta$ grob zu schätzen und daraus $S$ zu bestimmen, dann mit $S^{-1}$. Die zweistufige Schätzung aus Kapitel~3 ist der Sonderfall $W = (Z^{\top}Z)^{-1}$ --- optimal nur, wenn die Störungen gleich streuen.}
\end{minipage}
\end{center}

\schritt{Im Leitfall}
Zwei Instrumente für die Zuzahlung --- die Reform und ein regionaler Tarifunterschied:

\begin{center}
{\small
\renewcommand{\arraystretch}{1.48}
\begin{tabular}{@{}P{3.4cm}P{2.4cm}@{}}
\toprule
\textbf{Verfahren} & \textbf{je Euro}\\
\midrule
nur Reform & $-0{,}0980$\\
nur Region & $-0{,}0987$\\
zweistufig, beide & $-0{,}0982$\\
GMM, effizient & $\mathbf{-0{,}0982}$\\
\bottomrule
\end{tabular}}
\end{center}
\kursiv{Bei gleichmäßig streuenden Störungen fallen zweistufige und effiziente Schätzung zusammen --- wie der Satz es verlangt.}
\end{merksatz}

\section{Die überschüssigen Bedingungen prüfen}

\begin{merksatz}
\boxtitel{Satz 4.5 --- der Test von Hansen}
Gelten alle Momentenbedingungen, so ist
\[ J = n\,\bar g(\hat\theta)^{\top}\,\hat S^{-1}\,\bar g(\hat\theta) \]
asymptotisch $\chi^{2}$-verteilt mit so vielen Freiheitsgraden, wie es Bedingungen mehr als Parameter gibt.

\schritt{Der Gedanke}
Mit $m$ Bedingungen und $k$ Parametern kann man $k$ Bedingungen exakt erfüllen. Die übrigen $m-k$ müssen \kursiv{von selbst} nahe null sein --- und wie nahe, sagt ihre Streuung. $J$ misst den gewichteten Abstand der überschüssigen Momente von null. $\blacksquare$

\schritt{Im Leitfall}
\begin{center}
{\small
\renewcommand{\arraystretch}{1.48}
\begin{tabular}{@{}P{3.2cm}P{2.0cm}P{2.4cm}@{}}
\toprule
\textbf{Lage} & $J$ & \textbf{Urteil}\\
\midrule
beide gültig & $0{,}09$ & \textcolor{gruen}{unauffällig}\\
Region wirkt direkt & $\mathbf{1782}$ & \textcolor{akzent}{verworfen}\\
\bottomrule
\end{tabular}}
\end{center}
\kursiv{Grenze bei einem Freiheitsgrad: $3{,}84$.}

\begin{center}
\begin{minipage}{0.9\textwidth}
\textbf{Der Test ist der Überidentifikationstest aus Kapitel~3, nun als allgemeines Prinzip --- und er sagt weiterhin nicht, welche Bedingung verletzt ist.}
\end{minipage}
\end{center}
\end{merksatz}

\section{Der Bruchfall}

\begin{warnung}
\boxtitelr{Zu viele Momente}
\begin{center}
{\small
\renewcommand{\arraystretch}{1.52}
\begin{tabular}{@{}P{2.8cm}P{5.8cm}@{}}
\toprule
\textbf{Die Verlockung} & Jede Funktion eines gültigen Instruments ist wieder ein gültiges Instrument: $z$, $z^{2}$, $z^{3}$, $z\cdot x_2$ \ldots\\
\textbf{Die Folge} & Mit vielen Momenten wird $\hat S$ schlecht geschätzt, die Gewichtung unsinnig und der Schätzer verzerrt --- in Richtung der kleinsten Quadrate.\\
\textbf{Die Regel} & Die Zahl der Momente soll gegenüber $n$ klein bleiben; der Test aus Satz 4.5 verliert bei vielen Momenten seine Aussagekraft.\\
\bottomrule
\end{tabular}}
\end{center}

\begin{center}
\begin{minipage}{0.88\textwidth}
\kursiv{\textbf{Das ist das Problem der schwachen Instrumente aus Kapitel~3 in anderer Gestalt}: Viele Instrumente, von denen jedes wenig erklärt, wirken wie ein einziges schwaches. Die Literatur über Panelmodelle in Kapitel~8 kennt den Fall besonders gut, weil dort die Instrumente mit jeder Zeitperiode von selbst mehr werden.}
\end{minipage}
\end{center}
\end{warnung}

\begin{warnung}
\boxtitelr{Und was Momente nicht liefern}
\begin{center}
\begin{minipage}{0.88\textwidth}
\kursiv{Die Momentenmethode schätzt Parameter, keine Verteilungen. \textbf{Wer wissen will, wie viele Versicherte \kursiv{gar nicht} mehr zum Arzt gehen}, braucht die Verteilung der Besuchszahlen --- und die liefert nur ein vollständiges Modell. Das nächste Kapitel gibt eines an, und bezahlt dafür mit einer Verteilungsannahme.}
\end{minipage}
\end{center}
\end{warnung}

\section{Übungsaufgaben}

\begin{aufgabe}
\boxtitel{Aufgabe 4.1 --- Momentenschätzer von Hand}
Die Zahl der Arztbesuche sei poissonverteilt. Leiten Sie den Momentenschätzer für $\lambda$ aus zwei verschiedenen Momenten her.

\loesung
\schritt{Schritt 1 --- Erstes Moment}
\[ \mathbb{E}[y - \lambda] = 0 \quad\Longrightarrow\quad \hat\lambda_1 = \bar y \]

\schritt{Schritt 2 --- Zweites Moment}
Bei Poisson gilt $\operatorname{Var}[y] = \lambda$:
\[ \mathbb{E}\bigl[(y-\lambda)^{2} - \lambda\bigr] = 0 \quad\Longrightarrow\quad \hat\lambda_2 = \text{Stichprobenvarianz} \]

\schritt{Schritt 3 --- Die Zahlen}
$5000$ simulierte Versicherte mit $\lambda = 3{,}2$:
\begin{center}
{\small
\renewcommand{\arraystretch}{1.48}
\begin{tabular}{@{}P{2.6cm}P{2.4cm}@{}}
\toprule
\textbf{Moment} & $\hat\lambda$\\
\midrule
Mittelwert & $3{,}234$\\
Varianz & $3{,}290$\\
\bottomrule
\end{tabular}}
\end{center}

\schritt{Schritt 4 --- Die Deutung}
\begin{center}
\begin{minipage}{0.9\textwidth}
\textbf{Zwei Momente, ein Parameter --- die Lage von Satz 4.4. Der Mittelwert ist der bessere Schätzer, und GMM gewichtet ihn entsprechend stärker.}
\end{minipage}
\end{center}

\begin{center}
\begin{minipage}{0.88\textwidth}
\kursiv{\textbf{Der Unterschied der beiden Schätzungen ist zugleich ein Test}: Weichen Mittelwert und Varianz deutlich voneinander ab, so ist die Zahl nicht poissonverteilt --- das ist die Prüfgröße aus Band 25, Aufgabe 4.3, nun als Spezialfall von Satz 4.5. Bei Arztbesuchen fällt der Test regelmäßig durch: Die Varianz übersteigt den Mittelwert um ein Mehrfaches.}
\end{minipage}
\end{center}
\end{aufgabe}

\begin{aufgabe}
\boxtitel{Aufgabe 4.2 --- Die Gewichtsmatrix verstehen}
Zwei Momente mit Varianzen $1$ und $4$, unkorreliert. Wie gewichtet GMM?

\loesung
\schritt{Schritt 1 --- Die Matrix}
\[ S = \begin{pmatrix} 1 & 0\\ 0 & 4\end{pmatrix}, \qquad W^{*} = S^{-1} = \begin{pmatrix} 1 & 0\\ 0 & \tfrac14\end{pmatrix} \]

\schritt{Schritt 2 --- Das Kriterium}
\[ Q(\theta) = \bar g_1^{2} + \tfrac14\,\bar g_2^{2} \]

\begin{center}
\begin{minipage}{0.9\textwidth}
\textbf{Das streuende Moment zählt nur ein Viertel --- im Verhältnis der inversen Varianzen.}
\end{minipage}
\end{center}

\schritt{Schritt 3 --- Mit Korrelation}
Sind die Momente mit $\rho = 0{,}9$ korreliert, so wird
\[ S^{-1} = \frac{1}{4(1-0{,}81)}\begin{pmatrix} 4 & -1{,}8\\ -1{,}8 & 1\end{pmatrix} = \begin{pmatrix} 5{,}26 & -2{,}37\\ -2{,}37 & 1{,}32\end{pmatrix} \]

\begin{center}
\begin{minipage}{0.88\textwidth}
\kursiv{\textbf{Die Nebendiagonale ist negativ}: Zwei fast gleichlautende Bedingungen werden nicht doppelt, sondern nach ihrem \kursiv{Unterschied} gewichtet. Das ist der Grund, warum das Hinzufügen fast gleicher Instrumente nichts bringt --- und warum die Gewichtsmatrix bei vielen Momenten kaum schätzbar ist: Sie muss die Korrelationen aller Paare kennen.}
\end{minipage}
\end{center}
\end{aufgabe}

\begin{aufgabe}
\boxtitel{Aufgabe 4.3 --- Ein Modell mit Erwartungen}
Ein Haushalt plant seinen Konsum so, dass $\mathbb{E}_t[u'(c_{t+1})\,R_{t+1}] = \delta\,u'(c_t)$ gilt --- die Eulergleichung der Konsumtheorie. Wie schätzt man $\delta$ ohne Verteilungsannahme?

\loesung
\schritt{Schritt 1 --- Die Momentenbedingung}
\[ \mathbb{E}_t\left[\frac{u'(c_{t+1})}{u'(c_t)}\,R_{t+1} - \frac1\delta\right] = 0 \]
--- bedingt auf alles, was zur Zeit $t$ bekannt ist.

\schritt{Schritt 2 --- Instrumente}
Jede zur Zeit $t$ bekannte Größe $z_t$ darf mit dem Ausdruck multipliziert werden; nach der Turmregel aus Band 25, Satz 7.3, bleibt der unbedingte Erwartungswert null:
\[ \mathbb{E}\left[z_t\left(\frac{u'(c_{t+1})}{u'(c_t)}R_{t+1} - \frac1\delta\right)\right] = 0 \]

\schritt{Schritt 3 --- Was das leistet}
\begin{center}
\begin{minipage}{0.9\textwidth}
\textbf{Man schätzt einen Parameter des Verhaltens, ohne anzugeben, wie Einkommen, Zinsen oder Preise verteilt sind.}
\end{minipage}
\end{center}

\begin{center}
\begin{minipage}{0.88\textwidth}
\kursiv{\textbf{Das ist die Anwendung, für die GMM 1982 entwickelt wurde} --- \Person{Hansen} und \Person{Singleton} schätzten damit die Risikoscheu amerikanischer Haushalte. Das Ergebnis war ernüchternd: Die Bedingungen wurden verworfen, und zwar deutlich. Die Differenz zwischen beobachteten Aktienrenditen und dem, was die Eulergleichung erlaubt, heißt seither Aktienprämienrätsel und ist nicht gelöst.}
\end{minipage}
\end{center}

\schritt{Schritt 4 --- Die Verbindung}
\begin{center}
\begin{minipage}{0.88\textwidth}
\kursiv{Die Momentenbedingung ist eine Martingalbedingung: Der Ausdruck in der Klammer ist eine Martingaldifferenz. \textbf{Band 25 lieferte den Begriff, dieses Kapitel benutzt ihn als Schätzgleichung}, und Kapitel~11 wird ihn zur Prüfung der Markteffizienz heranziehen.}
\end{minipage}
\end{center}
\end{aufgabe}

\begin{geschichte}
\boxtitelg{Von Pearson zu Hansen}
\Person{Karl Pearson} führte die Momentenmethode 1894 ein --- um eine Mischung zweier Normalverteilungen an Krabbenmessungen anzupassen, die sein Kollege \Person{Weldon} aus Neapel mitgebracht hatte.

\begin{center}
\begin{minipage}{0.9\textwidth}
\kursiv{Er setzte die ersten fünf theoretischen Momente den empirischen gleich und löste ein Polynom neunten Grades. \textbf{\Person{Fisher} verdrängte das Verfahren ab 1922 mit der Maximum-Likelihood-Methode}, die bei bekannter Verteilung effizienter ist --- und die Momentenmethode galt sechzig Jahre lang als überholt.}
\end{minipage}
\end{center}

\schritt{Die Wiederkehr}
\begin{center}
\begin{minipage}{0.9\textwidth}
\kursiv{\Person{Lars Peter Hansen} verallgemeinerte 1982 den Ansatz auf beliebig viele Bedingungen und bewies Satz 4.4 und 4.5. \textbf{Der Anlass war die Konsumtheorie aus Aufgabe 4.3}: Ökonomische Modelle liefern Momentenbedingungen, keine Verteilungen, und Fishers Methode war dafür nicht zu gebrauchen. Hansen erhielt 2013 den Nobelpreis --- gemeinsam mit \Person{Fama} und \Person{Shiller}, deren Streit über Markteffizienz Kapitel~11 aufnimmt.}
\end{minipage}
\end{center}

\schritt{Und die Einordnung}
\begin{center}
\begin{minipage}{0.9\textwidth}
\kursiv{\textbf{Fast alle Schätzer dieses Bandes sind Sonderfälle von GMM} --- kleinste Quadrate, Instrumente, zweistufige Schätzung, und auch Maximum Likelihood, wenn man die Ableitung der Likelihood als Moment auffasst. Die Methode ist weniger ein Verfahren als eine Sprache, in der sich die anderen Verfahren vergleichen lassen.}
\end{minipage}
\end{center}
\end{geschichte}

\begin{forschung}
\boxtitelp{Momentenschätzung heute}
\begin{center}
{\small
\renewcommand{\arraystretch}{1.45}
\begin{tabular}{@{}P{3.2cm}P{5.0cm}@{}}
\toprule
\textbf{Empirische Likelihood} & Gewichte statt Gewichtsmatrix --- bessere Eigenschaften bei kleinen Stichproben.\\
\textbf{Viele Momente} & Regularisierung der Gewichtsmatrix, Auswahl von Instrumenten.\\
\textbf{Orthogonalisierte Momente} & Momente, die gegen Fehler in Hilfsgrößen unempfindlich sind.\\
\bottomrule
\end{tabular}}
\end{center}

\schritt{Zur dritten Zeile}
\begin{center}
\begin{minipage}{0.88\textwidth}
\kursiv{Schätzt man Kontrollgrößen mit Lernverfahren aus Band 22, so überträgt sich deren Fehler auf den Effekt --- außer man wählt die Momentenbedingung so, dass ihre Ableitung nach den Hilfsgrößen null ist. \textbf{Das ist der Kern des \enquote{doppelt robusten} maschinellen Lernens} von \Person{Chernozhukov} und Mitarbeitern aus dem Jahr 2018 --- und derselbe Gedanke wie die doppelte Absicherung aus Band 23, Kapitel~8, nun als Eigenschaft eines Moments formuliert.}
\end{minipage}
\end{center}
\end{forschung}

\begin{merksatz}
\boxtitel{Was Sie jetzt können}
\begin{center}
\renewcommand{\arraystretch}{1.5}
\begin{tabular}{@{}P{3.0cm}P{5.6cm}@{}}
\toprule
\textbf{Definition 4.1} & Schätzer als Nullstelle eines Mittelwerts\\
\textbf{Satz 4.2} & Konsistenz aus dem Gesetz der großen Zahlen\\
\textbf{Satz 4.4} & inverse Varianz gewichtet optimal\\
\textbf{Satz 4.5} & überschüssige Momente prüfen\\
\textbf{Der Bruchfall} & viele Momente wirken wie ein schwaches\\
\bottomrule
\end{tabular}
\end{center}

\begin{center}
\begin{minipage}{0.88\textwidth}
\kursiv{Alle bisherigen Zielgrößen waren stetig und unbeschränkt. Das nächste Kapitel behandelt Entscheidungen --- ja oder nein --- und Größen, die bei null abgeschnitten sind. Dort braucht man wieder eine Verteilung, und man erfährt, was sie kostet.}
\end{minipage}
\end{center}
\end{merksatz}
